\chap[Part V. The Working Technician]{Chapter 21. Commissioning, Testing and Fault-Finding}
Everything in this book ends up in two activities: proving that a system works before it is handed over, and finding out why it has stopped working. This chapter gives a method for each.

\sect{21.1 The testing lifecycle}
\defn{21.1}{FAT} Factory Acceptance Test: the system tested at the integrator's premises, with simulated I/O, before it goes to site.

\defn{21.2}{SAT} Site Acceptance Test: the installed system tested on site with real field devices, usually witnessed by the owner.

\defn{21.3}{ITP} Inspection and Test Plan: the document that lists every test, how to do it, what should happen, and records what did happen. Each test is witnessed and signed, and every failure goes on a defects list.

A typical ITP layout:
\begin{tbl}{P{0.12\textwidth}L}
Column & Contents\\\midrule
Test ID & A unique number\\
Control function & What is being tested: enable/disable, duty/assist start, alarm, interlock\\
Test procedure & Exactly what the tester does, step by step\\
Expected result & What must happen, written before the test\\
Actual result & What did happen, written during the test\\
Pass or fail & Judged against the expected result, not against what seems reasonable\\
Witness initial & The witness confirms they saw it\\
Defect list? & Every fail is logged, fixed and retested\\
\end{tbl}
The header records the date, the implementer, the witness, the subsystem under test and the environment (simulation, or live on the system).

\sect{21.2 What to test}
\begin{enumerate}
\item \textbf{Point-to-point (I/O) checks.} Every input and output, from the field device to the SCADA screen: operate the device or simulate the signal (Chapter 2), confirm the right point changes, with the right value, units and state text.
\item \textbf{Functional tests.} Each control function in each mode: enable and disable, auto and manual, duty and standby, start and stop delays, interlocks.
\item \textbf{Boundary tests.} Values exactly at a threshold, just inside and just outside. If a start condition is ‘beyond setpoint minus deadband’, a value exactly on the threshold must \emph{not} trigger it.
\item \textbf{Limit tests.} Each adjustable setpoint above its maximum and below its minimum must be rejected and revert to the last good value; a value in range must be accepted and retained.
\item \textbf{Alarm tests.} Each alarm raises, shows the right text and priority, clears, and appears in the alarm history (Chapter 10).
\item \textbf{Failure tests.} Lost communications show stale quality (Rule 3.1); a failed instrument shows bad quality and the logic does the safe thing; equipment that becomes unavailable is handled as designed.
\item \textbf{Redundancy tests.} Server failover and switch-back, network path failure, power supply failure (Section 9.5).
\end{enumerate}
\rul{21.1}{Ask the design question when a test reveals one} Tests often expose behaviour nobody specified: should a lost ‘available’ signal stop a running assist pump immediately, or only stop it starting again? Do not decide it yourself on the day. Record it and ask the process engineer.

\sect{21.3 A method for fault-finding}
\begin{enumerate}
\item \textbf{Define the symptom exactly.} Which point, which screen, since when, all the time or sometimes, one site or many.
\item \textbf{Find what changed.} Most faults follow a change: a patch, a password, a configuration edit, a restart, a new device.
\item \textbf{Split the chain.} Check halfway along the path from field to screen (Section 1.3). The fault is on one side.
\item \textbf{Prove each layer with an independent tool.} A loop calibrator for the signal, the RTU's own diagnostics, a packet capture for the protocol, a generic OPC UA or MQTT client for the server, the server's own logs.
\item \textbf{Change one thing at a time,} and put it back if it made no difference.
\item \textbf{Record what you found and did,} so the next person, possibly you in a year, does not start again.
\end{enumerate}
\law{21.1}{Half-splitting} A chain of N links with one fault is found in at most ⌈log₂ N⌉ checks by testing the middle each time, against up to N − 1 checking link by link.
\[ k = \lceil \log_2 N \rceil \]
A 16-link chain from transmitter to client display takes at most 4 checks, against up to 15.

\rul{21.2}{Read the log before guessing} Servers, drivers and clients log the exact reason a connection or point failed. A status such as BadSubscriptionIdInvalid or ‘certificate rejected’ is worth more than an hour of trying things.

\sect{21.4 Worked case: a bridge that would not deliver data}
\wex{21.1} An old system's data is to reach a new SCADA through an OPC DA-to-UA bridge (Definition 6.2). The faults met, in order, and how each layer was proved:
\begin{enumerate}
\item The bridge service would not start: the account lacked the ‘Log on as a service’ right. Granted; the service started (Section 6.1).
\item The bridge could not read the DA server: DCOM permissions. Corrected; the bridge's log showed its address space built.
\item The SCADA could browse and select a point, but the point showed a configuration error. A generic OPC UA test client read the same point with good quality: the bridge was proven, and the fault was in the SCADA client.
\item After the bridge was restarted, the SCADA's log showed BadSubscriptionIdInvalid: its old subscription no longer existed (Section 6.3). Disabling and re-enabling the subscription object created a new one.
\item With signing enabled, the server log showed the client's certificate in the rejected folder; it was trusted. The connection then failed because the session was activating as anonymous: the user identity had been set on the point, not on the server connection object (Section 16.3).
\end{enumerate}
Each step was settled by a log line or an independent tool, never by guessing. The service then ran under a personal login, which works until that person's password changes; moving it to a dedicated service account (Rule 6.1) finished the job.

\sect{21.5 Commissioning security}
Security features need the same proof as control functions. Four rules:
\begin{enumerate}
\item \textbf{Check what is running, not what was intended.} Read the live configuration the component actually obeys (its generated configuration file, its open ports), not the settings page or the design document.
\item \textbf{Attack it from the network side,} from a client PC, not from the server itself, where loopback addresses and local firewall rules hide what the network sees. Confirm that plain ports are closed, anonymous and wrong passwords are refused, publishing outside permissions is refused, and revoked or foreign certificates are refused.
\item \textbf{Prove changes survive a restart} (Law 17.9): disable an account, restart the service, and confirm it is still disabled.
\item \textbf{A test that cannot fail proves nothing.} Where possible, break the protection deliberately and confirm the test catches it. And have someone other than the author review the design: the author's own tests only find the faults the author imagined.
\end{enumerate}
\law{21.2}{Mutation score} Deliberately introduce faults (mutants) and count how many some test catches:
\[ \text{score} = \frac{\text{mutants detected}}{\text{mutants introduced}} \]
A score below 1 means some protection is unproven.

\inpractice{SecureLink} SecureLink ships a security report that reads the live installation and lists each check against its IEC 62443-4-2 requirement, and a separate test harness that attacks an installation from a client PC and files an HTML report. A sensible commissioning order for any secure broker: place it in the right zone with inbound ports open only from client subnets and all outbound blocked; verify the installer's hash; choose the bind address and confirm the server certificate names it; enrol two administrators with multi-factor sign-in; issue one certificate per device; set the narrowest topic permissions; forward the audit log off the machine; take and test an encrypted backup; resolve every failed check; then run the network-side attack test with real device certificates, including a revoked one.

\sect{21.6 Handover}
A system is not finished until the next person can maintain it: as-built drawings and point lists; configuration exported and committed (Section 13.5); passwords and certificates recorded in the proper vault with their expiry dates; the ITP signed with all defects closed or accepted; and the failover and backup procedures tested, not just written.

\sect{Exercises}
\begin{enumerate}
\item What is the difference between FAT and SAT? (Factory with simulation, against site with real devices)
\item A start condition is ‘DO below setpoint minus deadband’. What should a reading exactly on that threshold do? (Not trigger a start)
\item How many checks at most to find one fault in a 32-link chain by half-splitting? (5)
\item A generic OPC UA client reads good values but the SCADA shows a configuration error. Where is the fault? (In the SCADA's point or connection configuration)
\item Why run security tests from a client PC rather than the server? (Loopback and local firewall rules hide what the network sees)
\item Ten mutants are introduced and eight are caught. Find the score and what it means. (0.8: two protections are unproven)
\item Name four things a system needs before handover. (Section 21.6)
\end{enumerate}
